\documentclass{article}
\usepackage{spconf,amsmath,graphicx,amssymb,cite,booktabs,url}
\usepackage{multirow}
% \usepackage{tikz}
% \usetikzlibrary{positioning, shapes.geometric}
% \usetikzlibrary{calc}
 \usepackage{enumitem}
\setlist[itemize]{noitemsep} 
% Example definitions.
% ––––––––––
\def\x{{\mathbf x}}
\def\L{{\cal L}}

\newcommand{\eegtrans}{EEGXF}

% Title.
% ——
% \title{On the Interplay of Temporal Context and Model Architecture for Naturalistic EEG Decoding}
% \title{Temporal Context and Architecture: \\A Benchmark for Naturalistic EEG Decoding}
\title{TEMPORAL CONTEXT AND ARCHITECTURE:\\A BENCHMARK FOR NATURALISTIC EEG DECODING}

% Author.
% ------
\name{Mehmet Ergezer}
%\thanks{Dr. Ergezer holds concurrent appointments as an Associate Professor at Wentworth Institute of Technology and as an Amazon Visiting Academic. This paper describes work performed at Wentworth Institute of Technology and is not associated with Amazon.}}
\address{School of Computing \& Data Science, Wentworth Institute of Technology \\
Boston, MA USA \\
         Email: ergezerm@wit.edu}

\begin{document}

\maketitle

\begin{abstract}
We study how model architecture and temporal context interact in naturalistic EEG decoding. Using the HBN movie‑watching dataset, we benchmark five architectures, CNN, LSTM, a stabilized Transformer (\eegtrans{}), S4, and S5, on a 4-class task across segment lengths from 8 s to 128 s. 
% we benchmark CNN, LSTM, a stabilized Transformer (\eegtrans{}), and a structured state‑space model (S5) on a 4‑class task across segment lengths from 8s to 128s. 
Accuracy improves with longer context: at 64 s, S5 reaches 98.7\%$\pm$0.6 and CNN 98.3\%$\pm$0.3, while S5 uses $\sim$20$\times$ fewer parameters than CNN. To probe real‑world robustness, we evaluate zero‑shot cross‑frequency shifts, cross‑task OOD inputs, and leave‑one‑subject‑out generalization. S5 achieves stronger cross‑subject accuracy but makes over‑confident errors on OOD tasks; \eegtrans{} is more conservative and stable under frequency shifts, though less calibrated in‑distribution. These results reveal a practical efficiency–robustness trade‑off: S5 for parameter‑efficient peak accuracy; \eegtrans{} when robustness and conservative uncertainty are critical. 


% We study how model architecture and temporal context interact in naturalistic EEG decoding. Using the HBN movie‑watching dataset, we benchmark CNN, LSTM, a stabilized Transformer (\textbf{EEGXF}), and a structured state‑space model (\textbf{S5}) on a 4‑class task across segment lengths from 8s to 128s. Accuracy improves with longer context, with a CNN peaking at \textbf{97.2\%$\pm$0.9} and S5 at \textbf{95.2\%$\pm$1.7} for 64s segments, while S5 is \textbf{$\sim20\times$} more parameter‑efficient than CNN. To probe real‑world robustness, we evaluate zero‑shot cross‑frequency shifts, cross‑task OOD inputs, and leave‑one‑subject‑out (LOSO) generalization. S5 achieves stronger cross‑subject accuracy in LOSO, but tends to be overconfident on OOD tasks; EEGXF is more conservative and stable under frequency shifts. These results surface a practical \textbf{efficiency–robustness trade‑off}: S5 for efficient, high‑accuracy settings; EEGXF when robustness and calibrated uncertainty are critical. We release code, logs, and models for reproducibility.

%%%
% We investigate how model architecture and temporal context interact in naturalistic EEG decoding. Using the Healthy Brain Network dataset, we benchmark five architectures, CNN, LSTM, Transformer, \eegtrans{} (a modernized EEG Transformer), and S5 (a structured state-space model), on a 4-class movie classification task across segment lengths from 8s to 256s. While CNNs perform well on short segments, both S5 and CNNs excel at longer durations, though S5 is significantly more parameter-efficient and ultimately achieves the highest accuracy.
%S5 and CNNs excel at longer durations; though S5 remains more parameter-efficient.
%S5 consistently excels in accuracy and training efficiency at longer durations.


% To assess real-world generalization, we conduct zero-shot cross-task, cross-frequency, and leave-one-subject-out (LOSO) testing. These reveal not only significant performance drops, but also fundamental differences in how architectures handle uncertainty. S5 tends to make high-confidence errors on out-of-distribution data, while \eegtrans{} collapses to a neutral, low-confidence state, indicating superior uncertainty handling and robustness. Our results highlight a critical trade-off between S5's efficiency and \eegtrans's safer generalization, offering practical insights for deploying EEG models. We release all code and models for reproducibility.
\end{abstract}


%
\begin{keywords}
% EEG decoding, state-space models, SSM, S5, Transformers,  naturalistic stimuli
EEG, Neural decoding, State‑space models (S4/S5), Transformers, Naturalistic audio-visual stimuli
\end{keywords}
%

\section{Introduction}
\label{sec:intro}

Decoding brain activity from EEG during naturalistic tasks like movie-watching is a challenging problem, requiring models that capture long-range temporal dependencies from noisy, continuous recordings~\cite{shirazi2024hbn}. While deep learning has been widely applied to EEG, traditional architectures face fundamental trade-offs. CNNs are efficient but constrained by local receptive fields~\cite{roy2019deep, schirrmeister2017deep}, while Transformers offer global context via attention~\cite{vaswani2017attention}, leading to many adaptations for EEG~\cite{song2023eeg, wang2024dbconformer, abibullaev2023transformerreview}, but suffer from quadratic scaling in time and memory.

To address these limitations, Structured state-space Models (SSMs) like S4~\cite{gu2022efficientlymodelinglongsequences} and S5/Mamba~\cite{gu2023mamba} have emerged, offering a compelling balance of long-range modeling and near-linear computational complexity. Recent studies have applied SSMs to EEG foundation models~\cite{wang2025eegmamba}, drowsiness detection~\cite{siddhad2024drowzee}, and clinical tasks~\cite{tran2024eegssm, li2025dgmamba}. However, a systematic comparison of how these modern architectures perform against varying temporal context lengths has been missing.

In this work, we investigate the interplay between model architecture and temporal context. We hypothesize that sequence models like S5 and Transformers benefit more from longer EEG segments than local models like CNNs. %S5 and CNN exhibit the strongest gains with longer temporal context.
%Fig. 1 shows CNN and S5 both scale strongly, with CNN nearly matching S5 at 64 s (98.3% vs 98.7%)
To test this, we benchmark five architectures on a 4-class naturalistic decoding task using HBN data across segments from 8 s to 128 s, including a modernized Transformer baseline, \eegtrans{}, to ensure a robust comparison. To move beyond simple accuracy metrics, we introduce a suite of challenging generalization tests across unseen cognitive tasks, sampling frequencies, and subjects (leave-one-subject-out). This comprehensive evaluation allows us to characterize a fundamental trade-off: S5 offers state-of-the-art speed and accuracy but is prone to overconfident errors on out-of-distribution data, whereas \eegtrans{} provides greater robustness at a higher computational cost. While other notable architectures like EEG-Conformer ~\cite{song2023eeg} and foundation models like EEGMamba ~\cite{wang2025eegmamba} exist, our work focuses on a controlled comparison of fundamental architectural classes (CNN, LSTM, SSM, Transformer) under a unified training and evaluation protocol, a comparison which is currently missing in the literature.

% To test this, we benchmark five architectures on a 4-class naturalistic decoding task using HBN data across segments from 8s to 256s. To ensure a robust comparison, we also develop \eegtrans{}, a modernized Transformer baseline for EEG. Our contributions are: (1) a systematic evaluation of five architectures against temporal context length; (2) the introduction of \eegtrans{} as a strong baseline; (3) novel cross-frequency and leave-one-subject-out (LOSO) generalization tests; and (4) the characterization of a critical performance-robustness trade-off between S5 and \eegtrans.

% \paragraph*{Our Contributions}
% \begin{itemize}
%     \item We present the first systematic evaluation of five neural architectures across EEG segment lengths from 8s to 256s.
%     \item We introduce \eegtrans{}, a modernized and robust Transformer baseline for EEG decoding.
%     \item We evaluate generalization using three robust tests: (1) zero-shot cross-task inference on unseen cognitive states, (2) cross-frequency downsampling, and (3) leave-one-subject-out (LOSO) subject transfer.
%     \item We identify and characterize a fundamental trade-off: S5 offers state-of-the-art speed and accuracy but is prone to overconfident errors on out-of-distribution data, whereas \eegtrans{} provides greater robustness and safer uncertainty handling at a higher computational cost.
% \end{itemize}

% BibTeX entries to be added:
% - EEGMamba
% - DrowzEE-G-Mamba
% - EEG-SSM (Tran et al. 2024)
% - DG-Mamba (Li et al. 2025)

\section{Dataset and Preprocessing}
\label{sec:data}

% (Same as before, but optionally update to reflect 40+ subjects and multiple segment lengths)

We utilize EEG data from the Healthy Brain Network (HBN) initiative \cite{shirazi2024hbn}. Specifically, we use recordings from the movie-watching tasks, %where participants viewed "Despicable Me," "Diary of a Wimpy Kid," and "The Present."
where participants viewed ``Despicable Me,'' ``Diary of a Wimpy Kid,'' ``The Present,'' and a resting-state baseline.

The raw EEG data were downsampled to 250 Hz and the first 64 channels were selected for consistency. Our preprocessing pipeline consisted of:
\begin{enumerate}[itemsep=0cm,parsep=0.1cm]
    \item \textbf{Filtering:} Signals were first band-passed (1-50Hz) at 250Hz. For downsampling, an anti-aliasing low-pass filter was applied with a cutoff of 0.45 times the new target sampling rate before resampling. Segments were re-standardized after this process.%A 1-50 Hz bandpass filter.
    \item \textbf{Re-referencing:} Common average reference.
    \item \textbf{Segmentation:} The continuous EEG was segmented into $T$-second windows  with a 50\% overlap.
    \item \textbf{Normalization:} Each segment was standardized on a per-channel basis (zero mean, unit variance).
\end{enumerate}

Our benchmark uses data from 40 subjects, yielding a dataset of over 15,000 EEG segments. %The task is a 3-way classification to identify the movie from the EEG segment. 
The task is a 4-class classification to identify the stimulus: one of three movies or resting-state baseline.
The data was split into training (60\%), validation (20\%), and test (20\%) sets, stratified by subject and class.

\section{Model Architectures}
\label{sec:models}

% All models accept input tensors of shape (batch, 64,$T$) sequences; $T$ varies with segment length (e.g., T=2000 for 8s at 250Hz) and output 4-class logits.
We segmented recordings into $S\in\{8,16,32,64,128\}$‑second windows with 50\% overlap. All models accept inputs of shape (batch, 64, $T$), where $T$ depends on $S$ (e.g., $T=2000$ for 8 s at 250Hz).

\subsection{Structured state-space Model (S5)}
Our S5 classifier builds on recent advances in state-space modeling and is governed by a continuous-time linear state-space model:
\begin{align}
\mathbf{h}'(t) &= \mathbf{A} \mathbf{h}(t) + \mathbf{B} \mathbf{u}(t) \\
\mathbf{y}(t) &= \mathbf{C} \mathbf{h}(t) + \mathbf{D} \mathbf{u}(t)
\end{align}
where $\mathbf{u}(t)$ is the input, $\mathbf{h}(t)$ is the latent state, and $\mathbf{y}(t)$ is the output. In contrast to S4, which computes output via convolutional kernels in the frequency domain, S5 directly diagonalizes the state matrix $\mathbf{A}$ to enable efficient time-domain parallel scan computation.

% Each S5 layer includes a multi-input, multi-output (MIMO) SSM block, 
The typical S5 configuration uses 3 layers and a hidden dimension of 192 for cross-frequency experiments, consistent with the 182K parameter model shown in the main table.
% a learnable timescale vector, GELU nonlinearity, and LayerNorm. 
Our implementation consists of an input projection, a stack of 2--4 bidirectional S5 blocks, global average pooling, and a linear classifier. S5 models are trained using EMA and support irregular sampling and continuous-time generalization.

% We compare S5 against several widely used baselines:
% \begin{itemize}
%     \item \textbf{CNN}: A temporal CNN with three 1D convolutional blocks (kernel size 25, batch norm, ReLU, max pooling).
%     \item \textbf{LSTM}: A 2-layer bidirectional LSTM with a hidden size of 128 per direction.
%     \item \textbf{Transformer}: A 4-layer Transformer encoder with 8 attention heads and a model dimension of 256.
% \end{itemize}

%%%%
\subsection{\eegtrans{}: A Modernized and Stabilized Transformer Baseline} \label{sec:\eegtrans}

Standard Transformers often suffer from training instability on noisy EEG data. We observed that a vanilla baseline exhibited vanishing output variance, hindering learning. To create a robust baseline, we developed \eegtrans{}, an \textbf{EEG Transformer} that incorporates design principles from successful CNN and SSM architectures. Key modifications include a simplified input projection with ReLU activation, multiple BatchNorm layers for stabilization, a conservative encoder with \textit{norm\_first=True}, and high-variance weight initialization to encourage healthier gradient flow. This purpose-built architecture serves as our competitive Transformer baseline.

To aggregate temporal features, the \eegtrans{} replaces global average pooling with a learned attention pooling mechanism. A trainable query vector computes a weighted average over the sequence, allowing the model to prioritize more informative time steps.
A final linear classifier maps this embedding to the 4-class movie prediction.

% This architecture achieves 89.2\% accuracy on 8-second segments and is used as a competitive baseline throughout our experiments. While robust to sampling rate changes, it is 7.3$\times$ slower to train than S5 and less effective on longer segments.



\section{Experiments}
\label{sec:exp}

Models were trained for up to 100 epochs using the AdamW optimizer, a `ReduceLROnPlateau` scheduler, and early stopping. We used gradient clipping and EMA to stabilize training. %Results for the 64s segment are reported as mean~$\pm$~std over three random seeds; due to  computational cost, other results are from a single seed. 
Results for 32 s, 64 s, and 128 s are reported as mean~$\pm$~std over 3 seeds (CNN/\eegtrans{}/S5); 8s and 16s are single‑seed due to cost.
Reported training times are wall‑clock to early‑stopped convergence on a single NVIDIA RTX 3090 Ti. %, averaged over 3 seeds for 64s. 
Unless otherwise noted, Figures~\ref{fig:segcurve}-\ref{fig:efficiency_scatter} show single-seed results for five architectures (CNN, \eegtrans{}, LSTM, S4, S5), whereas Table~\ref{tab:main} reports mean $\pm$ std over three seeds for our focal models (CNN, \eegtrans{}, S5) at 32/64/128 s to conserve space.
% For reproducibility, all code, logs, and models will available  post-review\footnote{\url{https://anonymous.github.com/soon_to_be_linked}}.
All code and models will be released upon acceptance; a link to be added in camera‑ready.

% Models were trained for up to 100 epochs using the AdamW optimizer (LR $10^{-3}$, weight decay $10^{-4}$), a `ReduceLROnPlateau` scheduler, and early stopping. We used gradient clipping and EMA to stabilize training. Due to the extensive computational cost of the ablation, results are from a single seed; however, our robust generalization tests provide strong evidence for our findings. For reproducibility, all code, logs, and models are available at \url{https://github.com/TODO/eeg-s5-icassp}, using fixed random seeds on the public HBN dataset.

% \subsection{Baseline Models}
% We compare S5 against baselines including a temporal \textbf{CNN}, and a 2-layer bidirectional \textbf{LSTM}. %, and a standard 4-layer \textbf{Transformer}.

% All code and models are available at an anonymized link for review\footnote{\url{https://anonymous.4open.science/r/your-id}}.

% \section{Experiments}
% \label{sec:exp}

% Models were trained for up to 100 epochs using the AdamW optimizer (LR $10^{-3}$, weight decay $10^{-4}$), a ReduceLROnPlateau scheduler, and early stopping. We used gradient clipping and Exponential Moving Average (EMA) to stabilize training. For reproducibility, all code, logs, and models are available at \url{https://github.com/your-real-repo/eeg-s5-icassp}, using fixed random seeds on the public HBN dataset.
% %TODO UPDATE!!!

% Due to the extensive computational cost of the ablation study, results are reported from a single seed (42). While this is a limitation, our robust generalization tests (LOSO, cross-task, cross-frequency) provide strong evidence for our core findings.

% \subsection{Training Setup}
% Models were trained for up to 100 epochs (batch size 32) using the AdamW optimizer (LR $10^{-3}$, weight decay $10^{-4}$). We used a `ReduceLROnPlateau` scheduler and early stopping with a patience of 10 epochs. Gradient clipping (max norm 1.0) and Exponential Moving Average (EMA, decay 0.995) were used to stabilize training. While an exhaustive hyperparameter search was beyond the scope of this study, all models were trained with established learning rates and configurations to ensure a fair comparison.

% \subsection{Evaluation Metrics}
% We evaluated models on the test set using: Accuracy (\%), weighted F1-Score, number of trainable parameters, total training time (s), average inference time per sample (ms), and peak GPU memory usage (MB).

% \subsection*{Reproducibility}
% All source code and training logs are available at: \url{https://github.com/your-repo/s4-eeg-icassp}. The use of fixed random seeds and public data ensures full reproducibility.



\section{Results}
\label{sec:results}
Our primary results on subject-mixed splits are summarized in Figure~\ref{fig:segcurve} and Table~\ref{tab:main}\footnote{Subject-mixed: train/test split includes overlapping subjects. Subject-independent results are detailed in Sections 5.2, 5.3.}. 
Accuracy improves with longer context: at 64 s, S5 reaches $98.7\%\pm0.6$ and CNN $98.3\%\pm0.3$, 
while S5 uses $\sim$20$\times$ fewer parameters than CNN. 
CNN remains the fastest wall-clock at 64 s (9.8 $\pm$ 0.03 min), and S5 and \eegtrans{} have comparable training times (25.8 $\pm$ 6.2 vs.\ 25.1 $\pm$ 0.1 min). 

% Our primary results on subject-mixed splits are summarized in Figure 1 and Table 1. The experiment confirms our hypothesis: models with long-range capacity benefit more from temporal context. The CNN baseline is a remarkably strong performer, achieving a peak accuracy of 97.2\%~$\pm$~0.9 at 64s. S5 scales nearly as well, peaking at 95.2\%~$\pm$~1.7. %, but does so with over 20x fewer parameters and significantly faster training. Our engineered \eegtrans{} also scales effectively but peaks lower than the other two main contenders.
% S5 trains slightly faster than \eegtrans{}  and is far more parameter‑efficient than CNN; CNN remains the fastest wall‑clock at 64s.

The traditional LSTM and our simplified S4 baselines were not competitive. The LSTM failed to converge, and the S4 model, while reaching 73.7\% accuracy at 64 s, was prohibitively slow; details in Fig. 2.
%over 7 hours
Figure~\ref{fig:efficiency_scatter} further illustrates the accuracy-efficiency trade-offs, showing S5 and CNN as the most efficient high-performing architectures.

% Our primary results on subject-mixed splits are summarized in Figure~\ref{fig:segcurve} and Table~\ref{tab:results}. The experiment confirms our hypothesis: models with long-range capacity benefit more from temporal context. S5's accuracy scales best, rising from 55.3\% at 8s to a peak of 95.2\%~$\pm$~1.7 at 64s. The CNN baseline is also a strong performer, achieving 97.2\%~$\pm$~0.9 at 64s, while our engineered \eegtrans{} scales effectively but peaks lower.

% The traditional LSTM and our simplified S4 baselines were not competitive. The LSTM failed to converge, and the S4 model, while reaching 73.7\% accuracy at 64s, was prohibitively slow (over 7 hours vs. S5's $\sim$10 minutes). Figure~\ref{fig:efficiency_scatter} illustrates the accuracy-efficiency trade-offs, showing S5 and CNN as the most efficient high-performing architectures.

% Our primary results on subject-mixed splits are summarized in Figure~\ref{fig:segcurve} and Table~\ref{tab:results}. The experiment confirms our hypothesis: models with long-range capacity benefit more from temporal context. S5's accuracy scales best, rising from 55.3\% at 8s to a peak of 94.8\% at 64s, demonstrating its superior ability to leverage longer sequences. The CNN baseline is also a strong performer, achieving 90.3\% at 64s, while our engineered \eegtrans{} scales effectively to 75.3\% at 128s; its lower F1-score at this length suggests a tendency to over-predict dominant classes as context increases.

% The traditional LSTM and our simplified S4 baselines were not competitive. The LSTM failed to converge, achieving near-chance accuracy across all segment lengths. The S4 model, while showing improvement with context (peaking at 73.7\% at 64s), was prohibitively slow, taking over 7 hours to train compared to S5's 13 minutes, confirming S5's architectural and efficiency advantages. Figure~\ref{fig:efficiency_scatter} further illustrates the accuracy-efficiency trade-offs, showing S5 and CNN as the most efficient architectures.


\begin{figure}[t]
  \centering
    \includegraphics[width=0.95\linewidth]{fig1_accuracy_vs_segment_length.pdf}
  % \caption{Test accuracy versus segment length. S5 and CNN demonstrate the strongest performance scaling with temporal context. The x-axis is log-scaled.}
\caption{Test accuracy versus segment length (8–128 s). Single-seed comparison of five architectures (CNN, \eegtrans{}, LSTM, S4, S5).  S5 and CNN demonstrate the strongest performance scaling with temporal context; x-axis log-scaled.}
  %Test accuracy versus input segment length for EEG classification models. Our proposed S5 and \eegtrans{} architectures demonstrate superior performance across all temporal contexts, with S5 achieving optimal efficiency at longer segment lengths. The x-axis uses logarithmic scaling.}
  \label{fig:segcurve}
\end{figure}

% \begin{table}[t]
%   \centering
%   % \caption{Performance on 4-class subject-mixed EEG classification. For the 64s segment, results are mean~$\pm$~std over 3 seeds.}
%   \caption{Performance-efficiency trade-off at 64s segments (single seed). Accuracy is plotted against parameters (log-scale), with marker size scaled by training time. S5 achieves high accuracy with the fewest parameters and fastest training time, occupying the optimal top-left quadrant.}
%   \label{tab:results}
%   \resizebox{\columnwidth}{!}{%
%   \begin{tabular}{lrrrrr}
%     \toprule
%     \textbf{Model} & \textbf{Seg (s)} & \textbf{Acc. (\%)} & \textbf{F1} & \textbf{Time (min)} & \textbf{Params} \\
%     \midrule
%     \textbf{S5} & 8 & 55.3 & 0.553 & 3.9 & \textbf{183K} \\
%     (SSM) & 16 & 61.3 & 0.614 & 5.1 & \textbf{183K} \\
%      & 32 & 81.6 & 0.814 & 8.0 & \textbf{183K} \\
%      & 64 & 95.2 $\pm$ 1.7 & .954 $\pm$ .015 & 10.1 $\pm$ 1.0 & \textbf{183K} \\
%      & 128 & 92.2 & 0.895 & 9.6 & \textbf{183K} \\
%     \midrule
%     \eegtrans{} & 8 & 60.3 & 0.605 & 16.2 & 217K \\
%     (Transformer) & 16 & 63.4 & 0.635 & 22.9 & 217K \\
%      & 32 & 63.4 & 0.632 & 13.4 & 217K \\
%      & 64 & 84.2 $\pm$ 0.6 & .749 $\pm$ .010 & 11.6 $\pm$ 0.1 & 217K \\
%      & 128 & 75.3 & 0.593 & 9.9 & 217K \\
%     \midrule
%     CNN & 8 & 52.8 & 0.531 & \textbf{0.9} & 4.4M \\
%     (Local) & 16 & 64.7 & 0.650 & \textbf{1.2} & 4.4M \\
%      & 32 & 77.2 & 0.769 & \textbf{2.5} & 4.4M \\
%      & 64 & \textbf{97.2 $\pm$ 0.9} & \textbf{.953 $\pm$ .018} & \textbf{4.9 $\pm$ 0.6} & 4.4M \\
%      & 128 & 86.4 & 0.639 & \textbf{3.8} & 4.4M \\
%     \bottomrule
%   \end{tabular}
%   }
% \end{table}


% \begin{table}[t]
%   \centering
%   \caption{Performance on 4-class subject-mixed EEG classification. Results for 32s, 64s, and 128s segments are mean~$\pm$~std over 3 seeds.}
%   \label{tab:results}
%   \resizebox{\columnwidth}{!}{%
%   \begin{tabular}{lrrrrr}
%     \toprule
%     \textbf{Model} & \textbf{Seg (s)} & \textbf{Acc. (\%)} & \textbf{F1} & \textbf{Time (min)} & \textbf{Params} \\
%     \midrule
%     \textbf{S5} & 8 & 55.3 & 0.553 & 3.9 & \textbf{183K} \\
%     (SSM) & 16 & 61.3 & 0.614 & 5.1 & \textbf{183K} \\
%      & 32 & 94.4 $\pm$ 0.7 & .930 $\pm$ .013 & 19.9 $\pm$ 2.2 & \textbf{183K} \\
%      & 64 & \textbf{98.7 $\pm$ 0.6} & \textbf{.978 $\pm$ .011} & 25.8 $\pm$ 6.2 & \textbf{183K} \\
%      & 128 & \textbf{95.8 $\pm$ 0.6} & \textbf{.929 $\pm$ .007} & 17.8 $\pm$ 4.8 & \textbf{183K} \\
%     \midrule
%     \eegtrans{} & 8 & \textbf{60.3} & \textbf{0.605} & 16.2 & 217K \\
%     (Transformer) & 16 & 63.4 & 0.635 & 22.9 & 217K \\
%      & 32 & 80.1 $\pm$ 2.0 & .748 $\pm$ .029 & 35.8 $\pm$ 0.3 & 217K \\
%      & 64 & 81.8 $\pm$ 1.8 & .730 $\pm$ .019 & 25.1 $\pm$ 0.1 & 217K \\
%      & 128 & 76.7 $\pm$ 1.5 & .526 $\pm$ .032 & 6.3 $\pm$ 2.6 & 217K \\
%     \midrule
%     CNN & 8 & 52.8 & 0.531 & \textbf{0.9} & 4.4M \\
%     (Local) & 16 & \textbf{64.7} & \textbf{0.650} & \textbf{1.2} & 4.4M \\
%      & 32 & \textbf{94.5 $\pm$ 0.3} & \textbf{.931 $\pm$ .004} & \textbf{7.5 $\pm$ 0.02} & 4.4M \\
%      & 64 & 98.3 $\pm$ 0.3 & .976 $\pm$ .004 & \textbf{9.8 $\pm$ 0.03} & 4.4M \\
%      & 128 & 89.2 $\pm$ 1.5 & .684 $\pm$ .097 & \textbf{6.4 $\pm$ 0.06} & 4.4M \\
%     \bottomrule
%   \end{tabular}
%   }
% \end{table}
\begin{table}[t]
\centering
\caption{Performance on 4-class subject-mixed EEG classification.
32/64/128 s are mean$\pm$std over 3 seeds; 8/16s are single-seed.}
% Results for 32/64/128s are averaged over 3 seeds; 8s and 16s are single-seed.
\label{tab:main}
\resizebox{\columnwidth}{!}{%
\begin{tabular}{lcccc}
\toprule
Model & Seg (s) & Acc. (\%) & F1 & Time (min) \\
\midrule
\multirow{5}{*}{\shortstack{S5 \\ (SSM)}}
& 8   & 55.3 & 0.55 & 3.9 \\
& 16  & 61.3 & 0.61 & 5.1 \\
& 32  & $\mathbf{94.4}\pm0.7$ & $0.93\pm0.013$ & $19.9\pm2.2$ \\
& 64  & $\mathbf{98.7}\pm0.6$ & $0.98\pm0.011$ & $25.8\pm6.2$ \\
& 128 & $\mathbf{95.8}\pm0.6$ & $0.93\pm0.007$ & $17.8\pm4.8$ \\
\midrule
\multirow{5}{*}{\shortstack{EEGXF \\ (Transformer)}}
& 8   & 60.3 & 0.61 & 16.2 \\
& 16  & 63.4 & 0.64 & 22.9 \\
& 32  & $80.1\pm2.0$ & $0.75\pm0.029$ & $35.8\pm0.3$ \\
& 64  & $81.8\pm1.8$ & $0.73\pm0.019$ & $25.1\pm0.1$ \\
& 128 & $76.7\pm1.5$ & $0.53\pm0.032$ & $6.3\pm2.6$ \\
\midrule
\multirow{5}{*}{\shortstack{CNN \\ (Local)}}
& 8   & 52.8 & 0.53 & 0.9 \\
& 16  & 64.7 & 0.65 & 1.2 \\
& 32  & $94.5\pm0.3$ & $0.93\pm0.004$ & $7.5\pm0.02$ \\
& 64  & $98.3\pm0.3$ & $0.98\pm0.004$ & $9.8\pm0.03$ \\
& 128 & $89.2\pm1.5$ & $0.68\pm0.097$ & $6.4\pm0.06$ \\
\bottomrule
\end{tabular}
}
\end{table}


% \begin{table}[t]
%   \centering
%   \caption{Performance on 4-class subject-mixed EEG classification. S5 and CNN excel at longer contexts, while S5 remains the most parameter-efficient.}
%   \label{tab:results}
%   \resizebox{\columnwidth}{!}{%
%   \begin{tabular}{lrrrrr}
%     \toprule
%     \textbf{Model} & \textbf{Seg (s)} & \textbf{Acc. (\%)} & \textbf{F1} & \textbf{Time (min)} & \textbf{Params} \\
%     \midrule
%     \textbf{S5} & 8 & 55.3 & 0.553 & 3.9 & \textbf{183K} \\
%     (SSM) & 16 & 61.3 & 0.614 & 5.1 & \textbf{183K} \\
%      & 32 & 81.6 & 0.814 & 8.0 & \textbf{183K} \\
%      & 64 & \textbf{94.8} & \textbf{0.949} & 12.8 & \textbf{183K} \\
%      & 128 & \textbf{92.2} & \textbf{0.895} & 9.6 & \textbf{183K} \\
%     \midrule
%     EEGXF & 8 & 60.3 & 0.605 & 16.2 & 217K \\
%     (Transformer) & 16 & 63.4 & 0.635 & 22.9 & 217K \\
%      & 32 & 63.4 & 0.632 & 13.4 & 217K \\
%      & 64 & 73.1 & 0.731 & 24.1 & 217K \\
%      & 128 & 75.3 & 0.593 & 9.9 & 217K \\
%     \midrule
%     CNN & 8 & 52.8 & 0.531 & \textbf{0.9} & 4.4M \\
%     (Local Baseline) & 16 & 64.7 & 0.650 & \textbf{1.2} & 4.4M \\
%      & 32 & 77.2 & 0.769 & \textbf{2.5} & 4.4M \\
%      & 64 & 90.3 & 0.900 & \textbf{4.0} & 4.4M \\
%      & 128 & 86.4 & 0.639 & \textbf{3.8} & 4.4M \\
%     \bottomrule
%   \end{tabular}
%   }
% \end{table}


% ... after Table 1 and its footnotetext ...

To further analyze classification behavior, we examined the confusion patterns between the three movie classes at the 64 s segment length (Table~\ref{tab:confusion_summary}). Both S5 and CNN demonstrate high precision, misclassifying only 2.7\% of movie segments. In contrast, \eegtrans{} exhibits significant confusion between movies (32.4\%), explaining its lower F1-score and highlighting its difficulty in separating the nuanced features of the different naturalistic stimuli compared to S5 and CNN.

\begin{table}[h]
\centering
\caption{Inter-movie confusion analysis at 64 s segments. The metric shows the percentage of movie segments incorrectly classified as a different movie.}
\label{tab:confusion_summary}
% \resizebox{0.85\columnwidth}{!}{%
\begin{tabular}{lc}
\toprule
\textbf{Model} & \textbf{Movie Confusion Rate (\%) $\downarrow$} \\
\midrule
S5 & \textbf{2.7\%} \\
CNN & \textbf{2.7\%} \\
\eegtrans{} & 32.4\% \\
\bottomrule
\end{tabular}
% }
\end{table}

% ... Continue with the rest of the Results section ...

Figure~\ref{fig:efficiency_scatter} displays the performance vs efficiency across architectures. 
S5 models occupy the top-left region of the efficiency-accuracy space, showing they achieve strong performance with minimal resources. In contrast, \eegtrans{} models trade robustness for compute cost.


\begin{figure}[t]
  \centering
  \includegraphics[width=0.98\linewidth]{fig3_performance_tradeoff.pdf}
  % \caption{ Performance–efficiency trade‑off at 64s. Accuracy vs parameters (log‑scale); marker size $\underset{\sim}{\propto}$ mean training time (min). CNN is fastest, while S5 attains state‑of‑the‑art accuracy with \textbf{$\sim20\times$} fewer parameters than CNN; EEGXF is comparable in training time to S5 but less accurate.}
  \caption{Performance–efficiency trade-off at 64 s. Single-seed comparison of five architectures; accuracy vs.\ parameters (log-scale); marker size $\propto$ training time (min). CNN is fastest (4.4 M parameters) while S5 (0.2 M parameters) attains state-of-the-art accuracy with $\sim$20$\times$ fewer parameters than CNN; \eegtrans{} is comparable in training time but less accurate.}
  %\caption{Performance-efficiency trade-off analysis at 64-second segments. Test accuracy is plotted against model parameters (millions, log scale) with marker size proportional to training time. S5 achieves the optimal balance of high accuracy ($>90\%$) with minimal parameters ($<0.2M$), demonstrating superior efficiency compared to the larger CNN and slower \eegtrans{}}
  \label{fig:efficiency_scatter}
\end{figure}

% \subsection{Subject-Mixed vs. Subject-Independent Evaluation}
% A critical consideration in EEG decoding is the data splitting strategy. Our main results (Table~\ref{tab:results}) use a \textbf{subject-mixed} split, where segments from all subjects are pooled before being divided into training and test sets. While common, this may inflate performance by allowing models to learn subject-specific neural patterns.

% To provide a more realistic assessment of generalization, we also employ \textbf{subject-independent} splits in our cross-frequency and leave-one-subject-out (LOSO) evaluations. In these setups, the test set contains only subjects unseen during training. As we will show, this leads to a significant drop in accuracy (e.g., S5 performance falls from its subject-mixed peak of 94.8\% (at 64s) to 53.1\% in subject-independent tests (at 32s), highlighting the immense challenge of cross-subject generalization and the importance of reporting results under both conditions.


\subsection{Cross-Frequency Robustness}
To assess generalization across EEG sampling rates, we conducted a zero-shot evaluation on subject-independent splits. Models were trained at 250Hz and tested on signals downsampled to 128Hz and 64Hz using FIR filtering. The results are shown in Table~\ref{tab:crossfreq} and Figure~\ref{fig:crossfreq_bar}. S5's accuracy drops significantly (over 18 percentage points), highlighting a sensitivity to temporal resolution. In contrast, \eegtrans{} remains remarkably stable, with an accuracy drop of less than 2\%. This reveals a key trade-off: S5 offers superior peak performance and speed, but \eegtrans{} is preferable for applications with variable sampling rates.

% This is Table 2
\begin{table}[h]
\centering
\caption{Zero-shot accuracy at different sampling rates.}
\label{tab:crossfreq}
% \resizebox{0.8\columnwidth}{!}{%
\begin{tabular}{lccc}
\toprule
\textbf{Model} & \textbf{250Hz} & \textbf{128Hz} & \textbf{64Hz} \\
\midrule
S5 & 53.1\% & 37.5\% & 34.4\% \\
\eegtrans & 40.9\% & 39.7\% & 40.7\% \\
\bottomrule
\end{tabular}
% }
\end{table}

% Add this figure back in!
% \begin{figure}[h]
%   \centering
%   \includegraphics[width=0.9\linewidth]{frequency_robustness_comparison.pdf}
%   \caption{Zero-shot performance at different sampling frequencies. S5's accuracy degrades when tested on downsampled EEG, while \eegtrans{} remains robust.}
%   \label{fig:crossfreq_bar}
% \end{figure}

% \subsection{Cross-Frequency Robustness}
% To assess generalization across EEG sampling rates, we conducted a zero-shot evaluation of models trained at 256Hz and tested on signals downsampled to 128Hz and 64Hz using two proper anti-aliasing methods: time-domain FIR filtering and frequency-domain FFT truncation.

% \textbf{S5 Results:} S5 achieved 53.1\% accuracy at 256Hz but showed performance drops at lower frequencies. Accuracy decreased to 37.5\% at 128Hz and 34.4\% at 64Hz using FIR filtering, with nearly identical results under FFT-based downsampling. Despite high validation performance (99.6\%), S5 exhibited sensitivity to temporal resolution changes, highlighting the importance of fine-scale dynamics in its learned representations.

% \textbf{\eegtrans{} Results:} In this subject-independent setting, the improved \eegtrans{} reached 40.9\% test accuracy at 256Hz after better training (73.5\% validation). Surprisingly, its performance remained remarkably stable across frequencies: 39.7\% at 128Hz and 40.7\% at 64Hz (FIR). FFT-based results were similarly consistent. The \eegtrans{} showed minimal accuracy drop (\textless2\%), indicating strong robustness to sampling rate changes.

% While we do not report formal significance tests, the performance deltas observed are substantial. 
% Table~\ref{tab:crossfreq} compares FIR and FFT-based downsampling. Both methods yield near-identical performance, supporting the validity of the anti-aliasing pipeline.

% \textbf{Takeaways:} S5 offers superior peak performance and trains 12.5\texttimes{} faster than \eegtrans{} (12.6 min vs. 2.6 h), but is more sensitive to downsampling. \eegtrans, while slower and lower-performing overall, may be preferable in scenarios involving nonuniform or variable sampling rates.

% \begin{table}[h]
% \centering
% \caption{Zero-shot accuracy at different sampling rates using FIR and FFT downsampling.}
% \label{tab:crossfreq}
% \resizebox{\columnwidth}{!}{
% \begin{tabular}{lcccc}
% \toprule
% \textbf{Model} & \textbf{256Hz} & \textbf{128Hz (FIR / FFT)} & \textbf{64Hz (FIR / FFT)} \\
% \midrule
% S5 & 53.1\% & 37.5\% / 37.5\% & 34.4\% / 34.3\% \\
% \eegtrans & 40.9\% & 39.7\% / 40.2\% & 40.7\% / 41.4\% \\
% \bottomrule
% \end{tabular}
% }
% \end{table}

% \textbf{Takeaways:} S5 offers superior peak performance and trains over 12x faster than \eegtrans{} (12.6 min vs. 2.6 h), but is more sensitive to downsampling, with its accuracy dropping by over 18 percentage points. In contrast, \eegtrans{}, while slower, is remarkably stable, with an accuracy drop of less than 2\%, making it preferable for scenarios with variable sampling rates.

Figure~\ref{fig:crossfreq_bar} highlights the trade-off between speed and robustness: S5 trains faster and achieves higher peak accuracy, but \eegtrans{} is more stable under sampling frequency changes.

\begin{figure}[t]
  \centering
  \includegraphics[trim={0 0 0 1.8cm},clip, width=0.98\linewidth]{frequency_robustness_comparison.pdf}
  \caption{Zero-shot generalization performance at different sampling frequencies. 
  S5 shows a large drop in accuracy when tested on downsampled EEG, while \eegtrans{} 
  remains robust. This highlights a speed-accuracy-robustness trade-off across architectures.}
  \label{fig:crossfreq_bar}
\end{figure}


\subsection{Leave-One-Subject-Out Generalization}
To rigorously assess subject-independent generalization, we conducted a full leave-one-subject-out (LOSO) evaluation on 32 s segments. We ran 60 LOSO folds (multiple sessions per subject); significance tests used the 20 folds available for both models. The results (Table~\ref{tab:loso_results}) confirm \textbf{S5's architecture is superior for cross-subject generalization}, achieving a mean accuracy of 55.9\%. While these scores are substantially lower than subject-mixed results and exhibit high variance (accuracies per fold ranged from 21\% to 100\%), they highlight the profound difficulty of cross-subject EEG decoding. 

% \textbf{Statistical significance.} Across the 20 matched folds where both models were run, a paired t-test confirms S5 significantly outperforms \eegtrans{} in accuracy ($10.6 \pm 15.1$\,pp, $p < .01$) and F1-score ($0.25 \pm 0.16$, $p < 10^{-5}$). A McNemar test further confirms S5 is more likely to exceed 50\% accuracy per fold ($\chi^2 = 12.0$, $p < .001$). These findings reinforce that S5 generalizes better across subjects.

% To rigorously assess subject-independent generalization, we conducted a full leave-one-subject-out (LOSO) evaluation on 32s segments. In each fold, one subject was held out %entirely 
% for testing, while the models were trained on the remaining subjects. This prevents the model from learning subject-specific neural patterns and is a crucial test for real-world applicability.

% The results, summarized in Table~\ref{tab:loso_results}, confirm that \textbf{S5's architecture is superior for cross-subject generalization}. S5 achieves a mean accuracy of 55.9\%, significantly outperforming the lightweight \eegtrans{} baseline. While these scores are substantially lower than the subject-mixed results (cf. Table~\ref{tab:results}) and exhibit high variance across folds (from 21\% to 100\%), they highlight the profound difficulty of cross-subject EEG decoding. S5's ability to better capture generalizable temporal features, even in this challenging setting, underscores its strength as a sequence model for EEG. Due to the high computational cost of the Transformer architecture, the \eegtrans{} evaluation was limited to 20 folds.

% --- This table is based on your final JSON results ---
\begin{table}[h!]
\centering
\caption{Leave-One-Subject-Out (LOSO) performance on 32 s segments across 40 subjects. S5 demonstrates stronger cross-subject generalization.}
% \caption{Leave-One-Subject-Out (LOSO) performance on 32s segments. S5 demonstrates significantly stronger cross-subject generalization than the Transformer baseline.}
\label{tab:loso_results}
\resizebox{\linewidth}{!}{%
\begin{tabular}{lcccc}
\toprule
\textbf{Model} & \textbf{Mean Acc. (\%)} & \textbf{Std. Dev.} & \textbf{Folds Run} & \textbf{Mean Time/Fold} \\
\midrule
\textbf{S5} & \textbf{55.9} & 15.9 & 60 / 60 & \textasciitilde41 min \\
\eegtrans{} & 48.4 & 12.1 & 20 / 60 & \textasciitilde52 min \\
\bottomrule
\end{tabular}
}
\end{table}

% \subsection{Zero-Shot Cross-Task Generalization}
% To test how models behave on unseen cognitive states, we evaluated our movie-trained models on three out-of-distribution (OOD) active tasks. The results (Table~\ref{tab:cross_task}) reveal a striking difference in generalization strategy.

% \textbf{S5 exhibited overconfident specialization}, consistently misclassifying all active tasks as `Movie 3` with a high mean confidence of 60.0\%, indicating significant overconfidence. In stark contrast, \textbf{\eegtrans{} demonstrated a "conservative collapse"}, correctly identifying the OOD tasks as unfamiliar by defaulting to its most neutral known class (`Resting State`) with an appropriately low confidence of 26.0\%. This behavior is not a model failure but a sign of robust uncertainty quantification, a critical property for safe real-world deployment.

% % --- This table is the replacement for the figure ---
% \begin{table}[h]
% \centering
% \caption{Zero-shot cross-task generalization on three unseen active tasks. S5 shows overconfident misclassification, while EEGXF exhibits robust, low-confidence generalization to a neutral state.}
% % \caption{Cross-task generalization on unseen active tasks. S5 shows overconfident misclassification, while \eegtrans{} exhibits robust, low-confidence generalization to a neutral state.}
% \label{tab:cross_task}
% \resizebox{\linewidth}{!}{%
% \begin{tabular}{llcc}
% \toprule
% \textbf{Model} & \textbf{OOD Task Input} & \textbf{Dominant Prediction} & \textbf{Mean Confidence} \\
% \midrule
% \textbf{S5} & Symbol Search & Movie 3 & 60.0\% \\
% (Specialized) & Contrast Change & Movie 3 & 60.0\% \\
% & Spatial Memory & Movie 3 & 60.5\% \\
% \midrule
% \textbf{\eegtrans{}} & Symbol Search & Resting State & 26.0\% \\
% (Robust) & Contrast Change & Resting State & 26.0\% \\
% & Spatial Memory & Resting State & 26.0\% \\
% \bottomrule
% \end{tabular}
% }
% \end{table}


\textbf{Statistical significance.}
Across the 20 matched LOSO folds, S5 outperformed \eegtrans{} by $10.6 \pm 15.1$\,pp in accuracy ($t_{19} = 3.13$, $p = 0.0055$) and $0.25 \pm 0.16$ in F1-score ($t_{19} = 7.05$, $p = 1.0 \times 10^{-6}$), with medium-to-large effect sizes.
A McNemar test confirmed that S5 was significantly more likely to exceed 50\% accuracy across folds ($\chi^2 = 12.0$, $p = 4.9 \times 10^{-4}$).
These findings reinforce that S5 generalizes better than \eegtrans{} despite its overconfidence issues.
%Details and code will be released upon acceptance.%

% \subsubsection{Statistical comparison}
% To verify that the superior mean accuracy of S5 is not due to chance, we ran pair‑wise significance tests on the matched LOSO folds where both models were evaluated (20 folds).
% A paired t‑test shows that S5 outperforms the lightweight \eegtrans{} by 10.6 ± 15.1 pp in accuracy (t₍19₎ = 3.13, p = 0.006, Cohen’s d = 0.70 – medium effect) and by 0.25 ± 0.16 in F1‑score (t₍19₎ = 7.05, p ≈ 1 × 10⁻⁶, d = 1.58 – large effect).
% A McNemar test that asks “Which model exceeds 50 % chance level in a given fold?” yields χ² = 12.0, p = 4.9 × 10⁻⁴ (n₀₁ = 12, n₁₀ = 0), confirming that S5 is significantly more likely to beat chance than EEGXF.  
% These results corroborate the qualitative gap seen in Table 3 and reinforce the claim that S5 generalises better across unseen subjects despite its over‑confidence issues discussed in §5.4.

% \subsection{Zero-Shot Cross-Task Generalization}
% To test how models behave on unseen cognitive states, we analyzed out-of-distribution (OOD) performance. While statistical significance is not formally tested, the divergent behaviors between architectures are stark and qualitatively significant. 
% To test how models behave on unseen cognitive states, 
% We evaluated our movie-trained models on three out-of-distribution (OOD) active tasks, comprising a total of \textbf{N=245 segments} from Symbol Search, Contrast Change Detection, and Spatial Working Memory. As a control, we also evaluated performance on held-out segments from the in-distribution `Resting State` task. The results (Table~\ref{tab:cross_task}) reveal a striking difference in generalization strategy.

\subsection{Zero-Shot Cross-Task Generalization}
To test how models behave on unseen cognitive states, we evaluated our movie-trained models on three out-of-distribution (OOD) active tasks ($N=245$). The results reveal a striking difference in generalization strategy, supported by calibration metrics on the in-distribution test set (Table~\ref{tab:calibration}).

\textbf{S5 exhibited overconfident specialization.} While it is exceptionally well-calibrated on in-distribution data, achieving the best NLL and Brier and near‑best ECE (on par with CNN), it consistently misclassified all OOD tasks as `Movie 3' with a high mean confidence of 60.0\% (Table~\ref{tab:cross_task}). Its failure to express uncertainty on novel inputs is a critical limitation.

In stark contrast, \textbf{\eegtrans{} demonstrated a ``conservative collapse.''} Although less calibrated on in-distribution data, it correctly identified the OOD tasks as unfamiliar by defaulting to its most neutral known class (`Resting State') with an appropriately low confidence of 26.0\%. This behavior is not a failure but a sign of robust uncertainty quantification, a vital property for safe real-world deployment.

% \textbf{S5 exhibited overconfident specialization}, consistently misclassifying all active OOD tasks as `Movie 3` with a high 
% mean confidence of 60.0\%. In %stark 
% contrast, \textbf{\eegtrans{} demonstrated a ``conservative collapse"}, correctly identifying the OOD tasks as unfamiliar by defaulting to its most neutral known class (`Resting State`) with an appropriately low confidence of 26.0\%. This behavior is not a %model 
% failure but a sign of %robust 
% uncertainty quantification, a critical property for safe real-world deployment. Both models correctly classified the in-distribution `Resting State` control task with high accuracy.

\begin{table}[h]
\centering
\caption{Zero-shot cross-task generalization on unseen active tasks (N=245) and an in-distribution control task (N=525).}
\label{tab:cross_task}
\resizebox{\linewidth}{!}{%
\begin{tabular}{llcc}
\toprule
\textbf{Model} & \textbf{Task Input} & \textbf{Dominant Prediction} & \textbf{Mean Confidence} \\
\midrule
\multicolumn{4}{l}{\textit{Out-of-Distribution (Active Tasks)}} \\
\textbf{S5} & Symbol Search & Movie 3 & 60.0\% \\
(Specialized) & Contrast Change & Movie 3 & 60.0\% \\
& Spatial Memory & Movie 3 & 60.5\% \\
\cmidrule{2-4}
\textbf{\eegtrans{}} & Symbol Search & Resting State & 26.0\% \\
(Robust) & Contrast Change & Resting State & 26.0\% \\
& Spatial Memory & Resting State & 26.0\% \\
\midrule
\multicolumn{4}{l}{\textit{In-Distribution Control}} \\
S5 & Resting State & Resting State & $>95.0\%$ \\
\eegtrans{} & Resting State & Resting State & $>90.0\%$ \\
\bottomrule
\end{tabular}
}
\end{table}



% \begin{table}[h]
% \centering
% \caption{Calibration and uncertainty metrics for 64s models on the in-distribution test set. Lower values indicate better calibration.}
% \label{tab:calibration}
% \resizebox{\columnwidth}{!}{%
% \begin{tabular}{lccc}
% \toprule
% \textbf{Model} & \textbf{NLL $\downarrow$} & \textbf{Brier Score $\downarrow$} & \textbf{ECE (\%) $\downarrow$} \\
% \midrule
% S5 & \textbf{0.056 $\pm$ 0.008} & \textbf{0.006 $\pm$ 0.002} & \textbf{1.09 $\pm$ 0.12} \\
% \eegtrans{} & 0.999 $\pm$ 0.023 & 0.077 $\pm$ 0.005 & 13.41 $\pm$ 1.68 \\
% CNN & 0.079 $\pm$ 0.014 & 0.007 $\pm$ 0.002 & 0.96 $\pm$ 0.20 \\
% \bottomrule
% \end{tabular}
% }
% \end{table}

\begin{table}[h]
\centering
\caption{Calibration metrics (64 s, in-distribution). Lower is better.}
\label{tab:calibration}
\begin{tabular}{lccc}
\toprule
Model & NLL $\downarrow$ & Brier $\downarrow$ & ECE (\%) $\downarrow$ \\
\midrule
S5    & $\textbf{0.056} \pm 0.008$ & $\textbf{0.006} \pm 0.002$ & $1.09 \pm 0.12$ \\
EEGXF & $0.999 \pm 0.023$ & $0.077 \pm 0.005$ & $13.41 \pm 1.68$ \\
CNN   & $0.079 \pm 0.014$ & $0.007 \pm 0.002$ & $\textbf{0.96} \pm 0.20$ \\
\bottomrule
\end{tabular}
\end{table}

\section{Discussion and Conclusion}
\label{sec:conclusion}

Our systematic benchmark reveals a critical accuracy-efficiency-robustness trade-off in EEG decoding. In subject-mixed settings, both CNN and S5 architectures demonstrate powerful scaling with temporal context. However, S5 achieves this high performance with over 20x fewer parameters and greater training efficiency, establishing it as a leading architecture. Furthermore, in more challenging subject-independent tests (LOSO), S5's superiority in learning generalizable features becomes even more evident.

% However, this performance comes at the cost of robustness... % (The rest of the section is perfect as is)

% Our systematic benchmark reveals a critical accuracy-efficiency-robustness trade-off in EEG decoding. We show that S5's architecture is exceptionally suited for learning long-range temporal patterns, achieving the highest accuracy and a training speed over 10x faster within our benchmark than our robust Transformer baseline, \eegtrans.

However, this efficiency comes at the cost of robustness. Our novel cross-task generalization test demonstrates that S5 is prone to overconfident, incorrect predictions on unseen tasks. Conversely, \eegtrans's ``conservative collapse'' to a neutral, low-confidence state represents superior uncertainty handling, a vital property for any real-world BCI or clinical system where safety and reliability are paramount. This trade-off forces a practical choice: S5 is ideal for speed in controlled, subject-specific settings, while an architecture like \eegtrans{} is preferable for applications demanding robust generalization to novel cognitive states.

By characterizing this trade-off and releasing our code, we provide a foundation and a clear challenge for future research: to develop models that combine the efficiency of SSMs with the robust uncertainty awareness of Transformers.

\clearpage
\bibliographystyle{IEEEbib}
\bibliography{refs}

\end{document}
